% from main.tex
%%%%%%%%%%  scan idx 572  |  folio 559  %%%%%%%%%%
puissance. Par exemple sur le topos \((\Ens)\) annelé par le corps \(\mathbb Q\), la représentation \(\Gm\to\Gm\) définie par
\[
 r\rightsquigarrow|r|\qquad(r\in\mathbb Q^*)
\]
n'est pas une puissance. En particulier 6.2.3 ne reste pas valable pour n'importe quel topos localement annelé \(X\) et pour n'importe quelle représentation \(\alpha:G\to\Gm\).

\textbf{6.5.3.} En revanche 6.2.3 reste valable pour les représentations associées aux formules (6.2.8), (6.2.9) et (6.2.10) sur un topos annelé quelconque \(X=(\mathcal C,\mathcal O)\), parce que ces représentations sont «~algébriques~» dans le sens suivant :

Soient \(e\) l'objet final de \(X\), \(S=\Spec(\mathcal O(e))\), et considérons le foncteur
\[
 u:\mathcal C\longrightarrow(\Sch/S)
\]
qui associe à chaque \(Y\in\operatorname{Ob}(\mathcal C)\) le schéma affine \(\Spec(\mathcal O(Y))\) sur \(S\). Alors pour tout entier positif \(n\) le préfaisceau \(\operatorname{G\ell}(n)_X\) sur \(\mathcal C\) est composé avec \(u\) du préfaisceau correspondant \(\operatorname{G\ell}(n)_S\) sur \((\Sch/S)\). On dit qu'une représentation sur \(X\)
\begin{equation}\tag{6.5.4}
 \alpha:\operatorname{G\ell}(n)\longrightarrow\Gm
\end{equation}
est \emph{algébrique}, si elle provient par \(u\) d'une représentation
\begin{equation}\tag{6.5.5}
 \alpha_S:\operatorname{G\ell}(n)_S\longrightarrow\mathbb G_{m,S}
\end{equation}
de préfaisceaux en groupes sur \((\Sch/S)\), i.e., si l'on a
\[
 \alpha(s)=\alpha_S(s)
\]
pour tout \(Y\in\operatorname{Ob}(\mathcal C)\), \(s\in\operatorname{G\ell}(n)(Y)=\operatorname{G\ell}(n)_S(u(Y))\). La représentation \(\alpha\) dans (6.5.4) est bien entendu complètement déterminée par \(\alpha_S\). Il s'ensuit immédiatement que (6.2.3) reste valable dans un topos annelé quelconque, lorsque \(\alpha\) est algébrique. Donc en particulier les formules (6.2.8), (6.2.9) et (6.2.10) restent valables sur un topos annelé quelconque.

